\documentclass[12pt,a4paper,openany]{report}
\usepackage[UTF8]{ctex}
\usepackage{geometry}
\usepackage{hyperref}
\usepackage{listings}
\usepackage{xcolor}
\usepackage{graphicx}
\usepackage{booktabs}
\usepackage{longtable}
\usepackage{array}
\usepackage{enumitem}
\usepackage{fancyhdr}
\usepackage{titlesec}

% 页面设置
\geometry{left=2.5cm,right=2.5cm,top=2.5cm,bottom=2.5cm}
\raggedbottom % 避免页面拉伸产生的空白

% 颜色定义
\definecolor{codegreen}{rgb}{0,0.6,0}
\definecolor{codegray}{rgb}{0.5,0.5,0.5}
\definecolor{codepurple}{rgb}{0.58,0,0.82}
\definecolor{backcolour}{rgb}{0.95,0.95,0.92}

% 代码样式设置
\lstdefinestyle{mystyle}{
    backgroundcolor=\color{backcolour},
    commentstyle=\color{codegreen},
    keywordstyle=\color{magenta},
    numberstyle=\tiny\color{codegray},
    stringstyle=\color{codepurple},
    basicstyle=\ttfamily\footnotesize,
    breakatwhitespace=false,
    breaklines=true,
    captionpos=b,
    keepspaces=true,
    numbers=left,
    numbersep=5pt,
    showspaces=false,
    showstringspaces=false,
    showtabs=false,
    tabsize=2
}
\lstset{style=mystyle}

% 页眉页脚设置
\pagestyle{fancy}
\fancyhead[LE,RO]{\leftmark}
\fancyhead[RE,LO]{\thepage}
\fancyfoot[C]{\thepage}

% 章节标题设置
\titleformat{\chapter}[display]
{\normalfont\huge\bfseries\color{blue!70!black}}{\chaptertitlename\ \thechapter}{20pt}{\Huge}

% 超链接设置
\hypersetup{
    colorlinks=true,
    linkcolor=blue,
    filecolor=magenta,
    urlcolor=cyan,
    pdfpagemode=FullScreen,
}

% 文档信息
\title{\textbf{UDAKE 项目测试规范文档}}
\author{UDAKE 开发团队}
\date{\today}

\begin{document}

% 设置目录深度
\setcounter{tocdepth}{3}
\setcounter{secnumdepth}{3}

% 封面
\maketitle
\thispagestyle{empty}

% 目录
\tableofcontents

% 第一章：概述与测试策略
\chapter{概述与测试策略}

\section{文档目的}

本文档旨在为 UDAKE 项目提供全面的测试规范指南，涵盖前端、后端、通讯及功能模块的所有测试相关内容。本文档从测试角度出发，包含测试策略、测试范围、具体测试步骤以及所有功能的输入输出和API接口定义。

\section{项目概述}

UDAKE 是一个集成了实时插值、多目标优化、模型融合、路径规划等多个专业功能模块的空间数据分析与优化平台。项目采用前后端分离架构，前端使用 TypeScript + React/Vue，后端使用 Python FastAPI，支持 WebSocket 实时通讯。

\section{测试策略}

\subsection{测试分层策略}

\begin{enumerate}
    \item \textbf{单元测试}：测试独立的函数、类和组件
    \item \textbf{集成测试}：测试模块间的交互和数据流
    \item \textbf{系统测试}：测试完整的业务流程
    \item \textbf{性能测试}：测试系统性能指标和资源使用
    \item \textbf{安全测试}：测试安全防护措施和漏洞
    \item \textbf{E2E 测试}：测试完整的用户工作流
\end{enumerate}

\subsection{测试覆盖率要求}

\begin{table}[h]
\centering
\caption{测试覆盖率要求}
\begin{tabular}{|l|l|l|}
\hline
\textbf{覆盖率类型} & \textbf{目标值} & \textbf{说明} \\
\hline
语句覆盖率 & $\ge$ 80\% & 至少执行80\%的代码语句 \\
函数覆盖率 & $\ge$ 80\% & 至少调用80\%的函数 \\
分支覆盖率 & $\ge$ 80\% & 至少覆盖80\%的条件分支 \\
行覆盖率 & $\ge$ 80\% & 至少执行80\%的代码行 \\
\hline
\end{tabular}
\end{table}

\subsection{测试范围}

\subsubsection{前端测试范围}

\begin{itemize}
    \item 组件测试：UI组件的渲染和交互
    \item 服务层测试：API服务、WebSocket服务、配置服务等
    \item 状态管理测试：数据存储和状态更新
    \item 性能测试：加载时间、内存使用、Core Web Vitals
    \item 安全测试：XSS防护、CSP策略、敏感信息泄露
    \item E2E测试：完整用户流程和跨浏览器测试
\end{itemize}

\subsubsection{后端测试范围}

\begin{itemize}
    \item API测试：所有API端点的功能和性能
    \item 单元测试：核心算法和业务逻辑
    \item 集成测试：模块间交互和数据流
    \item 数据库测试：数据存储和查询
    \item 性能测试：响应时间、吞吐量、资源使用
    \item WebSocket测试：实时通讯和消息处理
\end{itemize}

\subsubsection{功能模块测试范围}

\begin{itemize}
    \item 实时插值系统：订阅管理、数据点添加、增量更新
    \item 多目标优化系统：优化算法、约束条件、帕累托解
    \item 模型融合系统：权重计算、融合策略、性能评估
    \item 路径规划系统：路径计算、约束验证、优化目标
    \item 采样建议系统：采样点生成、影响评估、优先级排序
    \item 不确定性分析：分级算法、风险计算、决策支持
\end{itemize}

\section{测试环境}

\subsection{开发环境}

\begin{itemize}
    \item 操作系统：macOS / Linux / Windows
    \item 前端框架：React/Vue + TypeScript
    \item 后端框架：Python 3.9+ + FastAPI
    \item 数据库：SQLite（开发）/ PostgreSQL（生产）
    \item 浏览器：Chrome / Firefox / Safari / Edge
\end{itemize}

\subsection{测试工具}

\begin{table}[h]
\centering
\caption{测试工具清单}
\begin{tabular}{|l|l|l|}
\hline
\textbf{工具类型} & \textbf{工具名称} & \textbf{用途} \\
\hline
前端单元测试 & Vitest & JavaScript/TypeScript单元测试 \\
前端E2E测试 & Playwright & 端到端测试和跨浏览器测试 \\
后端单元测试 & pytest & Python单元测试和集成测试 \\
代码覆盖率 & @vitest/coverage-v8 & 代码覆盖率统计 \\
API测试 & requests & HTTP API测试 \\
性能测试 & pytest-benchmark & 性能基准测试 \\
WebSocket测试 & pytest-asyncio & 异步WebSocket测试 \\
\hline
\end{tabular}
\end{table}

% 第二章：前端测试规范
\chapter{前端测试规范}

\section{测试框架}

\subsection{Vitest 框架}

Vitest 是项目使用的主要前端单元测试框架，提供快速的测试执行和现代化的测试API。

\subsubsection{配置文件}

\begin{lstlisting}[language=JavaScript, caption=vitest.config.js]
import { defineConfig } from 'vitest/config';

export default defineConfig({
  test: {
    environment: 'jsdom',
    globals: true,
    coverage: {
      provider: 'v8',
      reporter: ['text', 'json', 'html', 'lcov'],
      exclude: ['node_modules/', 'dist/', 'build/'],
      thresholds: {
        lines: 80,
        functions: 80,
        branches: 80,
        statements: 80
      }
    }
  }
});
\end{lstlisting}

\subsubsection{测试脚本}

\begin{lstlisting}[language=JSON, caption=package.json测试脚本]
{
  "scripts": {
    "test": "vitest --run",
    "test:watch": "vitest",
    "test:coverage": "vitest --run --coverage",
    "test:security": "vitest --run tests/security.test.ts",
    "test:e2e": "playwright test",
    "test:e2e:ui": "playwright test --ui",
    "test:performance": "vitest --run tests/performance-monitor.test.ts"
  }
}
\end{lstlisting}

\subsection{Playwright 框架}

Playwright 用于端到端测试，支持多浏览器和移动设备测试。

\subsubsection{配置文件}

\begin{lstlisting}[language=TypeScript, caption=playwright.config.ts]
import { defineConfig, devices } from '@playwright/test';

export default defineConfig({
  testDir: './tests/e2e',
  fullyParallel: true,
  retries: process.env.CI ? 2 : 0,
  workers: process.env.CI ? 1 : undefined,
  reporter: 'html',
  use: {
    baseURL: 'http://localhost:5173',
    trace: 'on-first-retry',
  },
  projects: [
    {
      name: 'chromium',
      use: { ...devices['Desktop Chrome'] },
    },
    {
      name: 'firefox',
      use: { ...devices['Desktop Firefox'] },
    },
    {
      name: 'webkit',
      use: { ...devices['Desktop Safari'] },
    },
    {
      name: 'Mobile Chrome',
      use: { ...devices['Pixel 5'] },
    },
    {
      name: 'Mobile Safari',
      use: { ...devices['iPhone 12'] },
    },
  ],
  webServer: {
    command: 'npm run dev',
    url: 'http://localhost:5173',
    reuseExistingServer: !process.env.CI,
  },
});
\end{lstlisting}

\section{测试文件结构}

\subsection{测试文件分类}

前端测试文件共39个，按功能分类如下：

\subsubsection{API相关测试（8个文件）}

\begin{itemize}
    \item \texttt{tests/api.test.js} - API基础功能测试
    \item \texttt{tests/api/api-service.test.ts} - API服务详细测试
    \item \texttt{tests/anomaly-detect-api.test.js} - 异常检测API测试
    \item \texttt{tests/config-api.test.js} - 配置API测试
    \item \texttt{tests/decision-thresholds-api.test.js} - 决策阈值API测试
    \item \texttt{tests/error-predict-api.test.js} - 误差预测API测试
    \item \texttt{tests/model-evaluation-api.test.js} - 模型评估API测试
    \item \texttt{tests/risk-calculate-api.test.js} - 风险计算API测试
    \item \texttt{tests/risk-report-api.test.js} - 风险报告API测试
    \item \texttt{tests/uncertainty-classify-api.test.js} - 不确定性分类API测试
\end{itemize}

\subsubsection{功能模块测试（12个文件）}

\begin{itemize}
    \item \texttt{tests/advancedfilter.test.js} - 高级过滤器测试
    \item \texttt{tests/batchoperations.test.js} - 批量操作测试
    \item \texttt{tests/confirmdialog.test.js} - 确认对话框测试
    \item \texttt{tests/datacomparison.test.js} - 数据对比测试
    \item \texttt{tests/exportenhancer.test.js} - 导出增强测试
    \item \texttt{tests/feedbackcollector.test.js} - 反馈收集器测试
    \item \texttt{tests/formvalidator.test.js} - 表单验证器测试
    \item \texttt{tests/historymanager.test.js} - 历史管理器测试
    \item \texttt{tests/offlinemanager.test.js} - 离线管理器测试
    \item \texttt{tests/preferencespanel.test.js} - 偏好设置面板测试
    \item \texttt{tests/templatedownloader.test.js} - 模板下载器测试
    \item \texttt{tests/thememanager.test.js} - 主题管理器测试
\end{itemize}

\subsubsection{性能和安全测试（5个文件）}

\begin{itemize}
    \item \texttt{tests/performance-monitor.test.ts} - 性能监控测试
    \item \texttt{tests/performancemonitor.test.js} - 性能监控器测试
    \item \texttt{tests/cache-performance.test.js} - 缓存性能测试
    \item \texttt{tests/security.test.ts} - 安全测试套件
    \item \texttt{tests/performance-benchmark.js} - 性能基准测试
\end{itemize}

\subsubsection{服务层测试（3个文件）}

\begin{itemize}
    \item \texttt{tests/websocket.test.ts} - WebSocket服务测试
    \item \texttt{tests/location-service.test.js} - 位置服务测试
    \item \texttt{tests/cache.test.js} - 缓存系统测试
\end{itemize}

\subsubsection{UI和交互测试（4个文件）}

\begin{itemize}
    \item \texttt{tests/gesture-enhancement.test.js} - 手势增强测试
    \item \texttt{tests/keyboardmanager.test.js} - 键盘管理器测试
    \item \texttt{tests/loadingmanager.test.js} - 加载管理器测试
    \item \texttt{tests/skeletonloader.test.js} - 骨架屏加载器测试
\end{itemize}

\subsubsection{E2E测试（1个文件）}

\begin{itemize}
    \item \texttt{tests/e2e/workflow.test.ts} - 完整工作流E2E测试
\end{itemize}

\section{核心测试用例}

\subsection{API服务测试}

\subsubsection{缓存机制测试}

\begin{lstlisting}[language=JavaScript, caption=缓存机制测试用例]
describe('缓存机制', () => {
  it('GET 请求应该被缓存', async () => {
    mockFetch.mockResolvedValueOnce({
      ok: true,
      json: () => Promise.resolve({ data: 'test' })
    });

    const url = 'http://localhost:8000/api/test';
    const result1 = await api.request(url);
    const result2 = await api.request(url);

    expect(mockFetch).toHaveBeenCalledTimes(1);
    expect(result1).toEqual({ data: 'test' });
    expect(result2).toEqual({ data: 'test' });
  });

  it('POST 请求不应该被缓存', async () => {
    mockFetch.mockResolvedValue({
      ok: true,
      json: () => Promise.resolve({ ok: true })
    });

    const url = 'http://localhost:8000/api/test';
    await api.request(url, { method: 'POST', body: '{}' });
    await api.request(url, { method: 'POST', body: '{}' });

    expect(mockFetch).toHaveBeenCalledTimes(2);
  });

  it('clearCache 应该清除所有缓存', async () => {
    mockFetch.mockResolvedValue({
      ok: true,
      json: () => Promise.resolve({ data: 1 })
    });

    await api.request('http://localhost:8000/api/test');
    api.clearCache();
    await api.request('http://localhost:8000/api/test');

    expect(mockFetch).toHaveBeenCalledTimes(2);
  });
});
\end{lstlisting}

\subsubsection{错误处理测试}

\begin{lstlisting}[language=JavaScript, caption=错误处理测试用例]
describe('错误处理', () => {
  it('HTTP 错误应该抛出异常', async () => {
    mockFetch.mockResolvedValueOnce({
      ok: false,
      status: 400,
      json: () => Promise.resolve({ detail: '服务器错误' })
    });
    await expect(api.request(url)).rejects.toThrow('服务器错误');
  });

  it('网络错误应该重试', async () => {
    mockFetch.mockRejectedValueOnce(new Error('Network error'));
    mockFetch.mockRejectedValueOnce(new Error('Network error'));
    mockFetch.mockResolvedValueOnce({
      ok: true,
      json: () => Promise.resolve({ data: 'success' })
    });

    const result = await api.get('/retry-test');
    expect(result).toEqual({ data: 'success' });
  });
});
\end{lstlisting}

\subsection{WebSocket服务测试}

\begin{lstlisting}[language=TypeScript, caption=WebSocket服务测试用例]
describe('WebSocketService', () => {
  let wsService: WebSocketService;

  beforeEach(() => {
    wsService = new WebSocketService('ws://localhost:8000');
  });

  it('应该能够连接到服务器', async () => {
    await wsService.connect('client-001');
    expect(wsService.isConnected).toBe(true);
  });

  it('应该能够发送消息', async () => {
    await wsService.connect('client-001');
    wsService.send({
      type: 'test',
      data: { message: 'hello' },
      timestamp: new Date().toISOString()
    });
  });

  it('应该能够接收消息', (done) => {
    wsService.on('test', (message) => {
      expect(message.data).toEqual({ message: 'hello' });
      done();
    });
  });

  it('连接失败时应该自动重连', async () => {
    const mockSocket = {
      connect: vi.fn().mockRejectedValue(new Error('Connection failed'))
    };
    // 测试重连逻辑
  });
});
\end{lstlisting}

\subsection{性能监控测试}

\begin{lstlisting}[language=TypeScript, caption=性能监控测试用例]
describe('性能监控', () => {
  let monitor: PerformanceMonitor;

  beforeEach(() => {
    monitor = new PerformanceMonitor();
  });

  it('应该能够获取 LCP', () => {
    const lcpEntry = { startTime: 1500 };
    mockPerformance.getEntriesByType.mockReturnValueOnce([lcpEntry]);
    const vitals = monitor.getWebVitals();
    expect(vitals.lcp).toBe(1500);
  });

  it('应该检测 LCP 过高', () => {
    mockPerformance.getEntriesByName.mockReturnValueOnce([{ startTime: 3000 }]);
    const result = monitor.checkPerformanceThresholds();
    expect(result.passed).toBe(false);
    expect(result.issues.some(issue => issue.includes('LCP'))).toBe(true);
  });

  it('应该监控内存使用', () => {
    mockPerformance.memory.usedJSHeapSize = 50 * 1024 * 1024; // 50MB
    mockPerformance.memory.totalJSHeapSize = 100 * 1024 * 1024; // 100MB
    const memoryUsage = monitor.getMemoryUsage();
    expect(memoryUsage.used).toBe(50 * 1024 * 1024);
  });
});
\end{lstlisting}

\subsection{安全测试}

\begin{lstlisting}[language=TypeScript, caption=安全测试用例]
describe('安全测试', () => {
  it('应正确转义用户输入', () => {
    const userInput = '<script>alert("XSS")</script>';
    const escapedInput = userInput
      .replace(/</g, '\&lt;')
      .replace(/>/g, '\&gt;');
    expect(escapedInput).not.toContain('<script>');
  });

  it('应验证文件路径，防止路径遍历攻击', () => {
    const validPath = '/home/user/data/file.txt';
    const invalidPath = '../../../etc/passwd';
    expect(isValidPath(validPath)).toBe(true);
    expect(isValidPath(invalidPath)).toBe(false);
  });

  it('应检查CSP策略', () => {
    const cspHeader = "default-src 'self'; script-src 'self' 'unsafe-inline'";
    expect(validateCSP(cspHeader)).toBe(true);
  });
});
\end{lstlisting}

\subsection{E2E测试}

\begin{lstlisting}[language=TypeScript, caption=E2E测试用例]
test('应该能够处理数据导入流程', async ({ page }) => {
  await page.goto('http://localhost:5173');
  const importButton = page.locator('button').filter({ hasText: /导入|import/i });
  await importButton.click();
  await expect(page.locator('input[type="file"]')).toBeVisible();
});

test('应该能够响应式布局', async ({ page }) => {
  await page.goto('http://localhost:5173');
  await page.setViewportSize({ width: 1920, height: 1080 });
  await page.setViewportSize({ width: 768, height: 1024 });
  await page.setViewportSize({ width: 375, height: 667 });
});

test('应该能够完成完整的插值工作流', async ({ page }) => {
  await page.goto('http://localhost:5173');
  // 1. 上传数据
  await page.locator('input[type="file"]').setInputFiles('test-data.geojson');
  // 2. 配置参数
  await page.locator('#variogram-model').selectOption('spherical');
  // 3. 启动插值
  await page.locator('button#start-kriging').click();
  // 4. 等待完成
  await expect(page.locator('.status-completed')).toBeVisible();
});
\end{lstlisting}

\section{测试执行}

\subsection{运行单元测试}

\begin{lstlisting}[language=bash, caption=运行单元测试]
# 运行所有测试
npm test

# 运行测试并监听变化
npm run test:watch

# 运行测试并生成覆盖率报告
npm run test:coverage

# 运行特定测试文件
npx vitest tests/api.test.js

# 运行匹配模式的测试
npx vitest --pattern "api"
\end{lstlisting}

\subsection{运行E2E测试}

\begin{lstlisting}[language=bash, caption=运行E2E测试]
# 运行所有E2E测试
npm run test:e2e

# 运行E2E测试并显示UI
npm run test:e2e:ui

# 运行特定浏览器
npx playwright test --project=chromium

# 调试模式
npx playwright test --debug
\end{lstlisting}

% 第三章：后端测试规范
\chapter{后端测试规范}

\section{测试框架}

\subsection{pytest 框架}

pytest 是项目使用的主要后端测试框架，支持单元测试、集成测试和异步测试。

\subsubsection{配置文件}

\begin{lstlisting}[language=Python, caption=pytest.ini]
[pytest]
testpaths = .
python_files = test_*.py
python_classes = Test*
python_functions = test_*
addopts =
    -v
    --tb=short
    --strict-markers
    --disable-warnings
    --cov=app
    --cov-report=term-missing
    --cov-report=html
asyncio_mode = auto
markers =
    slow: marks tests as slow
    integration: marks tests as integration tests
    unit: marks tests as unit tests
\end{lstlisting}

\subsubsection{测试依赖}

\begin{lstlisting}[language=text, caption=requirements.txt测试依赖]
pytest==7.4.3
pytest-asyncio==0.21.1
pytest-cov==4.1.0
pytest-benchmark==4.0.0
requests==2.31.0
httpx==0.25.2
\end{lstlisting}

\section{测试文件结构}

\subsection{测试文件分类}

后端测试文件共20个，按功能分类如下：

\subsubsection{核心模块测试（12个文件）}

\begin{itemize}
    \item \texttt{backend/tests/test\_model\_fusion.py} - 模型融合系统测试
    \item \texttt{backend/tests/test\_route\_planning.py} - 路径规划系统测试
    \item \texttt{backend/tests/test\_anomaly\_detector.py} - 异常检测器测试
    \item \texttt{backend/tests/test\_decision\_threshold\_analyzer.py} - 决策阈值分析器测试
    \item \texttt{backend/tests/test\_error\_predictor.py} - 误差预测器测试
    \item \texttt{backend/tests/test\_model\_evaluator.py} - 模型评估器测试
    \item \texttt{backend/tests/test\_risk\_index\_calculator.py} - 风险指数计算器测试
    \item \texttt{backend/tests/test\_sampling\_algorithms.py} - 采样算法测试
    \item \texttt{backend/tests/test\_spatial\_risk\_reporter.py} - 空间风险报告器测试
    \item \texttt{backend/tests/test\_uncertainty\_classifier.py} - 不确定性分类器测试
    \item \texttt{backend/tests/test\_websocket.py} - WebSocket服务测试
    \item \texttt{backend/test\_backend.py} - 后端集成测试
\end{itemize}

\subsubsection{实时插值系统测试（5个文件）}

\begin{itemize}
    \item \texttt{realtime\_interpolation/tests/test\_accuracy.py} - 准确性验证测试
    \item \texttt{realtime\_interpolation/tests/test\_performance.py} - 性能测试
    \item \texttt{realtime\_interpolation/tests/test\_integration.py} - 集成测试
    \item \texttt{realtime\_interpolation/tests/test\_cache\_system.py} - 缓存系统测试
    \item \texttt{realtime\_interpolation/tests/test\_incremental\_kriging.py} - 增量克里金测试
\end{itemize}

\subsubsection{多目标优化系统测试（1个文件）}

\begin{itemize}
    \item \texttt{multi\_objective\_optimization/tests/test\_core.py} - 核心算法测试
\end{itemize}

\subsubsection{API集成测试（2个文件）}

\begin{itemize}
    \item \texttt{api\_test.py} - 基础端点测试
    \item \texttt{api\_test\_complete.py} - 完整工作流测试
\end{itemize}

\section{核心测试用例}

\subsection{模型融合测试}

\begin{lstlisting}[language=Python, caption=权重计算器测试]
class TestWeightCalculator:
    """权重计算器测试"""

    @pytest.fixture
    def calculator(self):
        return WeightCalculator()

    @pytest.fixture
    def sample_metrics(self):
        return [
            ModelMetrics(
                model_id="model1",
                model_name="Model 1",
                rmse=1.2,
                mae=0.8,
                r2=0.85,
                stability=0.9
            ),
            ModelMetrics(
                model_id="model2",
                model_name="Model 2",
                rmse=1.5,
                mae=1.0,
                r2=0.80,
                stability=0.85
            ),
            ModelMetrics(
                model_id="model3",
                model_name="Model 3",
                rmse=1.0,
                mae=0.7,
                r2=0.90,
                stability=0.95
            )
        ]

    def test_calculate_weights_rmse_based(self, calculator, sample_metrics):
        """测试基于RMSE的权重计算"""
        weights = calculator.calculate_weights(
            sample_metrics,
            method='rmse_based'
        )
        assert len(weights) == 3
        assert all(0 <= w <= 1 for w in weights.values())
        assert abs(sum(weights.values()) - 1.0) < 1e-6

    def test_calculate_weights_r2_based(self, calculator, sample_metrics):
        """测试基于R²的权重计算"""
        weights = calculator.calculate_weights(
            sample_metrics,
            method='r2_based'
        )
        assert len(weights) == 3
        assert weights['model3'] > weights['model1']  # model3的R²最高
\end{lstlisting}

\subsection{WebSocket测试}

\begin{lstlisting}[language=Python, caption=WebSocket异步测试]
@pytest.mark.asyncio
async def test_websocket_connection():
    """测试 WebSocket 连接"""
    client = TestClient(app)
    with client.websocket_connect("/ws/test_client") as websocket:
        # 发送消息
        await websocket.send_json({
            "type": "subscribe_task",
            "task_id": "task_123"
        })

        # 接收确认消息
        response = await websocket.receive_json()
        assert response["type"] == "subscription_confirmed"
        assert response["task_id"] == "task_123"

@pytest.mark.asyncio
async def test_task_update_notification():
    """测试任务更新通知"""
    client = TestClient(app)
    with client.websocket_connect("/ws/test_client") as websocket:
        # 订阅任务
        await websocket.send_json({
            "type": "subscribe_task",
            "task_id": "task_123"
        })
        await websocket.receive_json()

        # 触发任务更新
        await notify_task_update("task_123", {"status": "running"})

        # 接收更新通知
        update = await websocket.receive_json()
        assert update["type"] == "task_update"
        assert update["data"]["status"] == "running"
\end{lstlisting}

\subsection{API测试}

\begin{lstlisting}[language=Python, caption=API端点测试]
def test_create_kriging_task():
    """测试创建克里金任务"""
    response = client.post("/start-kriging", json={
        "data_id": "data-123",
        "variogram_model": "spherical",
        "range": 100.0,
        "sill": 1.0,
        "nugget": 0.1,
        "grid_resolution": 50
    })
    assert response.status_code == 200
    data = response.json()
    assert "task_id" in data
    assert data["status"] == "pending"

def test_get_task_status():
    """测试获取任务状态"""
    # 先创建任务
    create_response = client.post("/start-kriging", json={...})
    task_id = create_response.json()["task_id"]

    # 查询状态
    response = client.get(f"/task-status/{task_id}")
    assert response.status_code == 200
    data = response.json()
    assert "status" in data
    assert "progress" in data

def test_get_prediction_result():
    """测试获取预测结果"""
    response = client.get(f"/result/prediction/{task_id}")
    assert response.status_code == 200
    data = response.json()
    assert "prediction" in data
    assert "variance" in data
\end{lstlisting}

\subsection{性能测试}

\begin{lstlisting}[language=Python, caption=性能基准测试]
def test_incremental_update_performance(benchmark):
    """测试增量更新性能"""
    # 准备测试数据
    subscription = create_test_subscription()
    new_point = DataPoint(x=50, y=50, value=10, id="test")

    # 基准测试
    result = benchmark(
        perform_incremental_update,
        subscription,
        new_point
    )

    # 性能断言
    assert result.update_time_ms < 100  # 更新时间小于100ms

def test_large_dataset_performance(benchmark):
    """测试大数据集性能"""
    # 创建大型方差网格（1000x1000）
    variance_grid = np.random.rand(1000, 1000)

    # 基准测试
    result = benchmark(
        optimize_sampling,
        variance_grid,
        n_samples=100
    )

    # 性能断言
    assert result.total_time < 10.0  # 总时间小于10秒
\end{lstlisting}

\section{测试执行}

\subsection{运行单元测试}

\begin{lstlisting}[language=bash, caption=运行后端单元测试]
# 运行所有测试
pytest

# 运行特定测试文件
pytest backend/tests/test_model_fusion.py

# 运行特定测试函数
pytest backend/tests/test_model_fusion.py::TestWeightCalculator::test_calculate_weights_rmse_based

# 运行标记的测试
pytest -m integration
pytest -m "not slow"

# 运行测试并生成覆盖率报告
pytest --cov=app --cov-report=html

# 详细输出
pytest -v
\end{lstlisting}

\subsection{运行异步测试}

\begin{lstlisting}[language=bash, caption=运行异步测试]
# 运行异步测试
pytest -m asyncio

# 运行特定异步测试
pytest backend/tests/test_websocket.py::test_websocket_connection
\end{lstlisting}

% 第四章：通讯与集成测试
\chapter{通讯与集成测试}

\section{通讯架构}

\subsection{通讯架构图}

\begin{verbatim}
前端                    后端
┌─────────────┐      ┌─────────────┐
│  API封装层   │ ───▶ │  FastAPI    │
│   Fetch     │      │  REST API   │
└─────────────┘      └──────┬──────┘
                            │
┌─────────────┐      ┌──────▼──────┐
│WebSocketSvc │ ◀───▶ │WebSocketSrv │
│ Socket.IO   │      │ FastAPI WS  │
└─────────────┘      └─────────────┘
                            │
                    ┌───────▼────────┐
                    │ 任务队列管理器  │
                    │ notify_task    │
                    └────────────────┘
\end{verbatim}

\section{WebSocket通讯测试}

\subsection{连接管理测试}

\begin{itemize}
    \item 测试连接建立和断开
    \item 测试自动重连机制
    \item 测试心跳检测
    \item 测试连接状态管理
\end{itemize}

\subsection{消息处理测试}

\begin{itemize}
    \item 测试消息发送和接收
    \item 测试消息确认机制
    \item 测试批量消息处理
    \item 测试消息队列管理
\end{itemize}

\subsection{任务订阅测试}

\begin{itemize}
    \item 测试任务订阅和取消订阅
    \item 测试任务更新通知
    \item 测试多任务并发订阅
    \item 测试订阅超时处理
\end{itemize}

\section{HTTP/REST API通讯测试}

\subsection{请求处理测试}

\begin{itemize}
    \item 测试GET请求缓存
    \item 测试POST请求不缓存
    \item 测试请求去重
    \item 测试错误重试机制
\end{itemize}

\subsection{响应处理测试}

\begin{itemize}
    \item 测试成功响应处理
    \item 测试错误响应处理
    \item 测试超时处理
    \item 测试响应数据验证
\end{itemize}

\section{集成测试}

\subsection{前后端集成测试}

\begin{lstlisting}[language=Python, caption=前后端集成测试]
def test_full_interpolation_workflow():
    """测试完整的插值工作流"""
    # 1. 前端上传数据
    upload_response = client.post("/upload-data", files={
        "file": ("test.geojson", open("test.geojson", "rb"), "application/json")
    })
    data_id = upload_response.json()["data_id"]

    # 2. 启动插值任务
    task_response = client.post("/start-kriging", json={
        "data_id": data_id,
        "variogram_model": "spherical",
        "range": 100.0,
        "sill": 1.0,
        "nugget": 0.1
    })
    task_id = task_response.json()["task_id"]

    # 3. 等待任务完成
    import time
    for _ in range(60):  # 最多等待60秒
        status_response = client.get(f"/task-status/{task_id}")
        status = status_response.json()["status"]
        if status == "completed":
            break
        time.sleep(1)

    # 4. 获取结果
    result_response = client.get(f"/result/prediction/{task_id}")
    assert result_response.status_code == 200
    result = result_response.json()
    assert "prediction" in result
\end{lstlisting}

\subsection{WebSocket与HTTP集成测试}

\begin{lstlisting}[language=Python, caption=WebSocket与HTTP集成测试]
@pytest.mark.asyncio
async def test_http_task_creation_with_websocket_updates():
    """测试HTTP创建任务，WebSocket接收更新"""
    client = TestClient(app)

    # 1. 建立WebSocket连接
    with client.websocket_connect("/ws/test_client") as websocket:
        # 2. 订阅任务
        await websocket.send_json({
            "type": "subscribe_task",
            "task_id": "task_123"
        })
        await websocket.receive_json()  # 确认消息

        # 3. 通过HTTP创建任务
        task_response = client.post("/start-kriging", json={...})
        task_id = task_response.json()["task_id"]

        # 4. 通过WebSocket接收更新
        update = await websocket.receive_json()
        assert update["type"] == "task_update"
        assert update["task_id"] == task_id
\end{lstlisting}

% 第五章：功能模块测试
\chapter{功能模块测试}

\section{实时插值系统}

\subsection{功能概述}

实时插值系统支持增量更新和实时预测，主要用于处理时序空间数据。

\subsection{输入输出格式}

\subsubsection{创建订阅请求}

\begin{lstlisting}[language=JSON, caption=创建订阅请求]
{
    "subscription_id": "sub_001",
    "data_type": "generic",
    "spatial_extent": {
        "min_x": 0.0,
        "max_x": 100.0,
        "min_y": 0.0,
        "max_y": 100.0
    },
    "update_frequency": 5,
    "interpolation_params": {
        "model_type": "spherical",
        "sill": 1.0,
        "range": 10.0,
        "nugget": 0.1
    },
    "notification_config": {}
}
\end{lstlisting}

\subsubsection{添加数据点请求}

\begin{lstlisting}[language=JSON, caption=添加数据点请求]
{
    "subscription_id": "sub_001",
    "point": {
        "id": "point_001",
        "x": 50.0,
        "y": 50.0,
        "value": 10.5,
        "timestamp": "2026-03-16T10:00:00Z",
        "metadata": {}
    }
}
\end{lstlisting}

\subsubsection{查询预测响应}

\begin{lstlisting}[language=JSON, caption=查询预测响应]
{
    "x": 50.0,
    "y": 50.0,
    "prediction": 10.5,
    "variance": 0.5,
    "confidence": 0.95,
    "version": 1,
    "timestamp": "2026-03-16T10:00:00Z"
}
\end{lstlisting}

\subsection{测试用例}

\subsubsection{订阅管理测试}

\begin{lstlisting}[language=Python, caption=订阅管理测试]
def test_create_subscription():
    """测试创建订阅"""
    response = client.post("/subscriptions", json={
        "subscription_id": "sub_001",
        "data_type": "sensor",
        "spatial_extent": {
            "min_x": 0.0,
            "max_x": 100.0,
            "min_y": 0.0,
            "max_y": 100.0
        },
        "update_frequency": 5
    })
    assert response.status_code == 201
    data = response.json()
    assert data["subscription_id"] == "sub_001"

def test_get_subscription():
    """测试获取订阅信息"""
    response = client.get("/subscriptions/sub_001")
    assert response.status_code == 200
    data = response.json()
    assert data["subscription_id"] == "sub_001"
    assert "spatial_extent" in data

def test_delete_subscription():
    """测试删除订阅"""
    response = client.delete("/subscriptions/sub_001")
    assert response.status_code == 200

    # 验证删除
    response = client.get("/subscriptions/sub_001")
    assert response.status_code == 404
\end{lstlisting}

\subsubsection{增量更新测试}

\begin{lstlisting}[language=Python, caption=增量更新测试]
def test_incremental_update():
    """测试增量更新"""
    # 1. 创建订阅并添加初始数据
    client.post("/subscriptions", json={...})
    client.post("/updates", json={
        "subscription_id": "sub_001",
        "point": {"id": "p1", "x": 50, "y": 50, "value": 10}
    })

    # 2. 添加新数据点
    response = client.post("/updates", json={
        "subscription_id": "sub_001",
        "point": {"id": "p2", "x": 60, "y": 60, "value": 12}
    })
    assert response.status_code == 200

    # 3. 查询预测值
    query_response = client.get(
        "/subscriptions/sub_001/query?x=55&y=55"
    )
    assert query_response.status_code == 200
    result = query_response.json()
    assert "prediction" in result
    assert "variance" in result
\end{lstlisting}

\subsubsection{性能测试}

\begin{lstlisting}[language=Python, caption=性能测试]
def test_update_performance(benchmark):
    """测试更新性能"""
    subscription = create_test_subscription()
    new_point = DataPoint(x=50, y=50, value=10, id="test")

    result = benchmark(
        perform_incremental_update,
        subscription,
        new_point
    )

    assert result.update_time_ms < 100
    assert result.affected_points > 0
\end{lstlisting}

\section{多目标优化系统}

\subsection{功能概述}

多目标优化系统使用NSGA-II算法，在多个目标（方差、成本、可达性）之间寻找最优解。

\subsection{输入输出格式}

\subsubsection{优化请求}

\begin{lstlisting}[language=JSON, caption=优化请求]
{
    "variance_grid": {
        "shape": [100, 100],
        "x_coords": [0, 1, 2, ..., 99],
        "y_coords": [0, 1, 2, ..., 99],
        "values": [[0.1, 0.2, ...], ...]
    },
    "existing_points": [
        {"x": 10.5, "y": 20.3},
        {"x": 15.2, "y": 25.7}
    ],
    "n_samples": 20,
    "weights": {
        "variance": 0.5,
        "cost": 0.3,
        "accessibility": 0.2
    },
    "constraints": {
        "boundary": {
            "type": "Polygon",
            "coordinates": [[[0, 0], [100, 0], [100, 100], [0, 100], [0, 0]]]
        },
        "min_distance": 50,
        "budget": 10000
    },
    "algorithm": "NSGA-II",
    "algorithm_params": {
        "population_size": 100,
        "n_generations": 100,
        "crossover_prob": 0.9,
        "mutation_prob": 0.1
    },
    "is_async": true
}
\end{lstlisting}

\subsubsection{优化结果}

\begin{lstlisting}[language=JSON, caption=优化结果]
{
    "success": true,
    "message": "获取成功",
    "data": {
        "task_id": "task_1234567890",
        "status": "completed",
        "pareto_solutions": [
            {
                "id": 1,
                "objectives": {
                    "variance": 0.125,
                    "cost": 8500,
                    "accessibility": -0.23
                },
                "sampling_points": [
                    {"id": 1, "x": 10.5, "y": 20.3, "variance": 0.15},
                    {"id": 2, "x": 15.2, "y": 25.7, "variance": 0.12}
                ],
                "metrics": {
                    "total_cost": 8500,
                    "avg_variance": 0.125,
                    "avg_accessibility": 0.23
                }
            }
        ],
        "recommended_solution": {
            "id": 5,
            "objectives": {
                "variance": 0.135,
                "cost": 7200,
                "accessibility": -0.28
            },
            "sampling_points": [...],
            "rank": 1,
            "crowding_distance": 0.45
        },
        "convergence_history": [...],
        "statistics": {
            "n_pareto_solutions": 15,
            "n_evaluations": 10000,
            "total_time": 8.2,
            "convergence_generation": 85
        }
    }
}
\end{lstlisting}

\subsection{测试用例}

\subsubsection{优化算法测试}

\begin{lstlisting}[language=Python, caption=优化算法测试]
def test_optimization_algorithm():
    """测试优化算法"""
    # 准备测试数据
    variance_grid = np.random.rand(50, 50)

    # 创建优化任务
    response = client.post("/multi-objective/optimize", json={
        "variance_grid": {
            "shape": [50, 50],
            "x_coords": np.arange(50).tolist(),
            "y_coords": np.arange(50).tolist(),
            "values": variance_grid.tolist()
        },
        "n_samples": 10,
        "weights": {
            "variance": 0.5,
            "cost": 0.3,
            "accessibility": 0.2
        }
    })
    assert response.status_code == 200
    task_id = response.json()["data"]["task_id"]

    # 等待完成
    wait_for_task_completion(task_id)

    # 获取结果
    result_response = client.get(f"/multi-objective/tasks/{task_id}/results")
    assert result_response.status_code == 200
    result = result_response.json()

    # 验证结果
    assert len(result["data"]["pareto_solutions"]) > 0
    assert "recommended_solution" in result["data"]
\end{lstlisting}

\subsubsection{约束条件测试}

\begin{lstlisting}[language=Python, caption=约束条件测试]
def test_boundary_constraint():
    """测试边界约束"""
    response = client.post("/multi-objective/optimize", json={
        "variance_grid": {...},
        "n_samples": 10,
        "constraints": {
            "boundary": {
                "type": "Polygon",
                "coordinates": [[[0, 0], [100, 0], [100, 100], [0, 100], [0, 0]]]
            }
        }
    })
    # 验证所有采样点都在边界内
    result = get_optimization_result(response.json()["task_id"])
    for solution in result["pareto_solutions"]:
        for point in solution["sampling_points"]:
            assert 0 <= point["x"] <= 100
            assert 0 <= point["y"] <= 100

def test_distance_constraint():
    """测试距离约束"""
    response = client.post("/multi-objective/optimize", json={
        "variance_grid": {...},
        "n_samples": 10,
        "constraints": {
            "min_distance": 50.0
        }
    })
    # 验证采样点之间的最小距离
    result = get_optimization_result(response.json()["task_id"])
    for solution in result["pareto_solutions"]:
        points = solution["sampling_points"]
        for i in range(len(points)):
            for j in range(i + 1, len(points)):
                distance = calculate_distance(points[i], points[j])
                assert distance >= 50.0

def test_budget_constraint():
    """测试预算约束"""
    response = client.post("/multi-objective/optimize", json={
        "variance_grid": {...},
        "n_samples": 10,
        "constraints": {
            "budget": 10000.0
        }
    })
    # 验证总成本不超过预算
    result = get_optimization_result(response.json()["task_id"])
    for solution in result["pareto_solutions"]:
        assert solution["metrics"]["total_cost"] <= 10000.0
\end{lstlisting}

\section{模型融合系统}

\subsection{功能概述}

模型融合系统支持多种融合策略，用于整合多个模型的预测结果。

\subsection{输入输出格式}

\subsubsection{创建融合任务请求}

\begin{lstlisting}[language=JSON, caption=创建融合任务请求]
{
    "models": [
        {
            "model_id": "model_1",
            "model_name": "普通克里金",
            "predictions": [10.5, 11.2, 10.8],
            "variances": [0.5, 0.6, 0.4]
        },
        {
            "model_id": "model_2",
            "model_name": "泛克里金",
            "predictions": [10.8, 11.0, 10.9],
            "variances": [0.4, 0.5, 0.3]
        }
    ],
    "config": {
        "strategy": "weighted_average",
        "weight_method": "rmse_based",
        "min_weight": 0.0,
        "max_weight": 1.0,
        "normalize": true,
        "smoothing": false,
        "enable_cross_validation": true,
        "enable_stability_check": true,
        "n_folds": 5
    },
    "true_values": [10.6, 11.0, 10.9]
}
\end{lstlisting}

\subsubsection{融合结果}

\begin{lstlisting}[language=JSON, caption=融合结果]
{
    "success": true,
    "result": {
        "task_id": "fusion_task_001",
        "fused_predictions": [10.65, 11.1, 10.85],
        "fused_variances": [0.45, 0.55, 0.35],
        "weights": {
            "model_1": 0.5,
            "model_2": 0.5
        },
        "metrics": {
            "rmse": 0.05,
            "mae": 0.04,
            "r2": 0.98
        },
        "strategy": "weighted_average",
        "weight_method": "rmse_based"
    }
}
\end{lstlisting}

\subsection{测试用例}

\subsubsection{权重计算测试}

\begin{lstlisting}[language=Python, caption=权重计算测试]
def test_weight_calculation():
    """测试权重计算"""
    response = client.post("/fusion/create-task", json={
        "models": [
            {
                "model_id": "model_1",
                "predictions": [10.5, 11.2, 10.8],
                "variances": [0.5, 0.6, 0.4]
            },
            {
                "model_id": "model_2",
                "predictions": [10.8, 11.0, 10.9],
                "variances": [0.4, 0.5, 0.3]
            }
        ],
        "config": {
            "strategy": "weighted_average",
            "weight_method": "rmse_based"
        },
        "true_values": [10.6, 11.0, 10.9]
    })
    assert response.status_code == 200
    result = response.json()

    # 验证权重
    weights = result["result"]["weights"]
    assert abs(sum(weights.values()) - 1.0) < 1e-6
    assert all(0 <= w <= 1 for w in weights.values())
\end{lstlisting}

\subsubsection{融合策略测试}

\begin{lstlisting}[language=Python, caption=融合策略测试]
def test_fusion_strategies():
    """测试不同融合策略"""
    strategies = ["weighted_average", "minimum_variance", "ensemble"]

    for strategy in strategies:
        response = client.post("/fusion/create-task", json={
            "models": [...],
            "config": {
                "strategy": strategy,
                "weight_method": "rmse_based"
            }
        })
        assert response.status_code == 200
        result = response.json()
        assert result["result"]["strategy"] == strategy
\end{lstlisting}

\section{路径规划系统}

\subsection{功能概述}

路径规划系统用于优化采样点的访问路径，支持多种优化目标和约束条件。

\subsection{输入输出格式}

\subsubsection{路径规划请求}

\begin{lstlisting}[language=JSON, caption=路径规划请求]
{
    "sampling_points": [
        {
            "id": "point1",
            "name": "采样点1",
            "latitude": 39.9042,
            "longitude": 116.4074,
            "priority": 5,
            "service_time": 10,
            "is_reachable": true
        }
    ],
    "start_point": {
        "id": "start",
        "name": "起点",
        "latitude": 39.8942,
        "longitude": 116.3974,
        "priority": 1,
        "service_time": 0,
        "is_reachable": true
    },
    "constraints": {
        "vehicle_type": "car",
        "time_windows": false,
        "priority_constraint": false,
        "max_distance": 100000,
        "max_duration": 3600,
        "max_cost": 1000
    },
    "optimization_goal": "shortest_distance",
    "algorithm": "tsp",
    "return_multiple_routes": true
}
\end{lstlisting}

\subsubsection{路径规划结果}

\begin{lstlisting}[language=JSON, caption=路径规划结果]
{
    "success": true,
    "routes": [
        {
            "route_id": "uuid",
            "point_sequence": ["start", "point1", "point2"],
            "segments": [
                {
                    "from": "start",
                    "to": "point1",
                    "distance": 500,
                    "duration": 60,
                    "cost": 10
                }
            ],
            "total_distance": 5000,
            "total_duration": 600,
            "total_cost": 100,
            "start_time": "2024-01-01T10:00:00",
            "end_time": "2024-01-01T10:10:00"
        }
    ],
    "best_route": {...},
    "statistics": {
        "total_routes": 1,
        "best_distance": 5000,
        "computation_time": 1.23
    },
    "warnings": []
}
\end{lstlisting}

\subsection{测试用例}

\subsubsection{路径计算测试}

\begin{lstlisting}[language=Python, caption=路径计算测试]
def test_route_calculation():
    """测试路径计算"""
    response = client.post("/api/route-planning/plan", json={
        "sampling_points": [
            {
                "id": "point1",
                "latitude": 39.9042,
                "longitude": 116.4074
            },
            {
                "id": "point2",
                "latitude": 39.9142,
                "longitude": 116.4174
            }
        ],
        "start_point": {
            "id": "start",
            "latitude": 39.8942,
            "longitude": 116.3974
        },
        "optimization_goal": "shortest_distance",
        "algorithm": "tsp"
    })
    assert response.status_code == 200
    result = response.json()

    # 验证路径
    assert len(result["routes"]) > 0
    assert "best_route" in result
    assert result["best_route"]["total_distance"] > 0
\end{lstlisting}

\subsubsection{约束验证测试}

\begin{lstlisting}[language=Python, caption=约束验证测试]
def test_constraint_validation():
    """测试约束验证"""
    response = client.post("/api/route-planning/plan", json={
        "sampling_points": [...],
        "constraints": {
            "max_distance": 1000,
            "max_duration": 600
        }
    })
    assert response.status_code == 200
    result = response.json()

    # 验证约束
    for route in result["routes"]:
        assert route["total_distance"] <= 1000
        assert route["total_duration"] <= 600
\end{lstlisting}

\section{其他功能模块}

\subsection{采样建议系统}

\subsubsection{功能概述}

采样建议系统根据方差栅格生成最优采样点推荐。

\subsubsection{测试用例}

\begin{lstlisting}[language=Python, caption=采样建议测试]
def test_sampling_recommendations():
    """测试采样建议生成"""
    response = client.get(
        "/api/sampling-recommendations/generate?task_id=task_001&strategy=hybrid&n_recommendations=20"
    )
    assert response.status_code == 200
    result = response.json()

    # 验证结果
    assert len(result["recommendations"]) == 20
    assert all("x" in rec and "y" in rec for rec in result["recommendations"])
    assert all("variance" in rec for rec in result["recommendations"])
\end{lstlisting}

\subsection{不确定性分级系统}

\subsubsection{功能概述}

不确定性分级系统根据预测方差对空间区域进行分级。

\subsubsection{测试用例}

\begin{lstlisting}[language=Python, caption=不确定性分级测试]
def test_uncertainty_classification():
    """测试不确定性分级"""
    response = client.post("/uncertainty/classify", json={
        "task_id": "task-001",
        "prediction": [[10.5, 11.2, 10.8], [11.0, 10.9, 11.3]],
        "variance": [[0.5, 0.8, 0.6], [0.7, 0.9, 0.4]],
        "x_coords": [120.1, 120.2, 120.3],
        "y_coords": [30.1, 30.2]
    })
    assert response.status_code == 200
    result = response.json()

    # 验证分级结果
    assert "statistics" in result
    assert "color_map" in result
    assert "critical_zones" in result
\end{lstlisting}

% 第六章：API接口测试
\chapter{API接口测试}

\section{API接口总览}

\subsection{接口分类统计}

\begin{table}[h]
\centering
\caption{API接口分类统计}
\begin{tabular}{|l|l|l|}
\hline
\textbf{分类} & \textbf{接口数量} & \textbf{主要功能} \\
\hline
核心插值 & 8个 & 插值任务、结果查询、报告生成 \\
采样优化 & 7个 & 采样建议、影响评估、预览 \\
不确定性分析 & 4个 & 分级、风险指数、决策阈值 \\
批量处理 & 11个 & 批量插值、报告、对比 \\
AI扩展 & 8个 & 模型评估、误差预测、异常检测、模型融合 \\
系统管理 & 25个 & 配置、任务队列、资源监控、性能报告 \\
专业系统 & 35个 & 实时插值、多目标优化、路径规划 \\
数据管理 & 18个 & 上传、导出、项目管理 \\
\hline
\textbf{总计} & \textbf{116个} & \\
\hline
\end{tabular}
\end{table}

\section{核心插值接口}

\subsection{接口列表}

\begin{longtable}{|l|l|l|l|}
\hline
\textbf{方法} & \textbf{路径} & \textbf{功能} & \textbf{文件位置} \\
\hline
\endhead
POST & \texttt{/start-kriging} & 启动克里金插值任务 & \texttt{插值任务接口.py:19} \\
GET & \texttt{/task-status/\{task\_id\}} & 查询任务状态 & \texttt{任务状态接口.py:17} \\
GET & \texttt{/result/prediction/\{task\_id\}} & 获取预测结果 & \texttt{结果查询接口.py:20} \\
GET & \texttt{/result/variance/\{task\_id\}} & 获取方差结果 & \texttt{结果查询接口.py:32} \\
GET & \texttt{/result/report/\{task\_id\}} & 生成克里金分析报告 & \texttt{报告生成接口.py:19} \\
GET & \texttt{/result/download/\{task\_id\}/\{filename\}} & 下载导出文件 & \texttt{结果查询接口.py:46} \\
GET & \texttt{/progress-detail/\{task\_id\}} & 获取任务进度详情 & \texttt{进度详情接口.py:17} \\
GET & \texttt{/estimated-remaining-time/\{task\_id\}} & 获取预计剩余时间 & \texttt{进度详情接口.py:29} \\
\hline
\end{longtable}

\subsection{接口测试用例}

\subsubsection{启动克里金任务}

\begin{lstlisting}[language=Python, caption=启动克里金任务测试]
def test_start_kriging():
    """测试启动克里金任务"""
    response = client.post("/start-kriging", json={
        "data_id": "data-123",
        "variogram_model": "spherical",
        "range": 100.0,
        "sill": 1.0,
        "nugget": 0.1,
        "grid_resolution": 50,
        "nlags": 6
    })
    assert response.status_code == 200
    data = response.json()
    assert "task_id" in data
    assert data["status"] == "pending"

def test_start_kriging_invalid_parameters():
    """测试启动克里金任务（无效参数）"""
    response = client.post("/start-kriging", json={
        "data_id": "data-123",
        "variogram_model": "invalid",
        "range": -100.0,
        "sill": -1.0
    })
    assert response.status_code == 422
\end{lstlisting}

\subsubsection{查询任务状态}

\begin{lstlisting}[language=Python, caption=查询任务状态测试]
def test_get_task_status():
    """测试查询任务状态"""
    # 先创建任务
    create_response = client.post("/start-kriging", json={...})
    task_id = create_response.json()["task_id"]

    # 查询状态
    response = client.get(f"/task-status/{task_id}")
    assert response.status_code == 200
    data = response.json()
    assert "status" in data
    assert "progress" in data
    assert data["task_id"] == task_id

def test_get_task_status_not_found():
    """测试查询任务状态（任务不存在）"""
    response = client.get("/task-status/nonexistent-task")
    assert response.status_code == 404
\end{lstlisting}

\section{实时插值系统接口}

\subsection{接口列表}

\begin{longtable}{|l|l|l|l|}
\hline
\textbf{方法} & \textbf{路径} & \textbf{功能} & \textbf{文件位置} \\
\hline
\endhead
POST & \texttt{/subscriptions} & 创建订阅 & \texttt{fastapi\_routes.py:64} \\
GET & \texttt{/subscriptions/\{subscription\_id\}} & 获取订阅信息 & \texttt{fastapi\_routes.py:92} \\
GET & \texttt{/subscriptions} & 列出所有订阅 & \texttt{fastapi\_routes.py:125} \\
DELETE & \texttt{/subscriptions/\{subscription\_id\}} & 删除订阅 & \texttt{fastapi\_routes.py:146} \\
POST & \texttt{/updates} & 添加数据点并触发增量更新 & \texttt{fastapi\_routes.py:167} \\
POST & \texttt{/subscriptions/\{subscription\_id\}/data-points} & 添加数据点 & \texttt{fastapi\_routes.py:200} \\
GET & \texttt{/subscriptions/\{subscription\_id\}/query} & 查询指定点的预测值 & \texttt{fastapi\_routes.py:233} \\
GET & \texttt{/subscriptions/\{subscription\_id\}/prediction} & 获取预测值 & \texttt{fastapi\_routes.py:257} \\
GET & \texttt{/cache/statistics} & 获取缓存统计信息 & \texttt{fastapi\_routes.py:281} \\
DELETE & \texttt{/cache} & 清除缓存 & \texttt{fastapi\_routes.py:311} \\
GET & \texttt{/system/status} & 获取系统状态 & \texttt{fastapi\_routes.py:333} \\
GET & \texttt{/system/metrics} & 获取性能指标 & \texttt{fastapi\_routes.py:364} \\
GET & \texttt{/system/alerts} & 获取告警列表 & \texttt{fastapi\_routes.py:388} \\
\hline
\end{longtable}

\subsection{接口测试用例}

\subsubsection{创建订阅}

\begin{lstlisting}[language=Python, caption=创建订阅测试]
def test_create_subscription():
    """测试创建订阅"""
    response = client.post("/subscriptions", json={
        "subscription_id": "sub_001",
        "data_type": "sensor",
        "spatial_extent": {
            "min_x": 0.0,
            "max_x": 100.0,
            "min_y": 0.0,
            "max_y": 100.0
        },
        "update_frequency": 5,
        "interpolation_params": {
            "model_type": "spherical",
            "sill": 1.0,
            "range": 10.0,
            "nugget": 0.1
        }
    })
    assert response.status_code == 201
    data = response.json()
    assert data["subscription_id"] == "sub_001"

def test_create_subscription_invalid_extent():
    """测试创建订阅（无效空间范围）"""
    response = client.post("/subscriptions", json={
        "subscription_id": "sub_001",
        "spatial_extent": {
            "min_x": 100.0,
            "max_x": 0.0,  # 无效：min > max
            "min_y": 0.0,
            "max_y": 100.0
        }
    })
    assert response.status_code == 422
\end{lstlisting}

\section{多目标优化系统接口}

\subsection{接口列表}

\begin{longtable}{|l|l|l|l|}
\hline
\textbf{方法} & \textbf{路径} & \textbf{功能} & \textbf{文件位置} \\
\hline
\endhead
POST & \texttt{/multi-objective/optimize} & 创建并执行优化任务 & \texttt{fastapi\_routes.py:77} \\
GET & \texttt{/multi-objective/tasks/\{task\_id\}/status} & 查询优化任务状态 & \texttt{fastapi\_routes.py:124} \\
GET & \texttt{/multi-objective/tasks/\{task\_id\}} & 获取优化任务信息 & \texttt{fastapi\_routes.py:152} \\
GET & \texttt{/multi-objective/tasks/\{task\_id\}/results} & 获取优化任务结果 & \texttt{fastapi\_routes.py:184} \\
DELETE & \texttt{/multi-objective/tasks/\{task\_id\}} & 取消优化任务 & \texttt{fastapi\_routes.py:211} \\
GET & \texttt{/multi-objective/config} & 获取系统配置 & \texttt{fastapi\_routes.py:239} \\
GET & \texttt{/multi-objective/tasks} & 获取优化任务列表 & \texttt{fastapi\_routes.py:295} \\
GET & \texttt{/multi-objective/tasks/\{task\_id\}/export/geojson} & 导出优化结果为GeoJSON & \texttt{fastapi\_routes.py:329} \\
GET & \texttt{/multi-objective/status} & 获取系统状态 & \texttt{fastapi\_routes.py:368} \\
\hline
\end{longtable}

\subsection{接口测试用例}

\subsubsection{创建优化任务}

\begin{lstlisting}[language=Python, caption=创建优化任务测试]
def test_create_optimization_task():
    """测试创建优化任务"""
    variance_grid = np.random.rand(50, 50)

    response = client.post("/multi-objective/optimize", json={
        "variance_grid": {
            "shape": [50, 50],
            "x_coords": np.arange(50).tolist(),
            "y_coords": np.arange(50).tolist(),
            "values": variance_grid.tolist()
        },
        "n_samples": 10,
        "weights": {
            "variance": 0.5,
            "cost": 0.3,
            "accessibility": 0.2
        },
        "constraints": {
            "boundary": {
                "type": "Polygon",
                "coordinates": [[[0, 0], [50, 0], [50, 50], [0, 50], [0, 0]]]
            },
            "min_distance": 5,
            "budget": 1000
        },
        "algorithm": "NSGA-II",
        "algorithm_params": {
            "population_size": 50,
            "n_generations": 50,
            "crossover_prob": 0.9,
            "mutation_prob": 0.1
        }
    })
    assert response.status_code == 200
    data = response.json()
    assert "task_id" in data["data"]
\end{lstlisting}

\section{模型融合接口}

\subsection{接口列表}

\begin{longtable}{|l|l|l|l|}
\hline
\textbf{方法} & \textbf{路径} & \textbf{功能} & \textbf{文件位置} \\
\hline
\endhead
POST & \texttt{/fusion/create-task} & 创建模型融合任务 & \texttt{模型融合接口.py:57} \\
GET & \texttt{/fusion/task/\{task\_id\}/status} & 获取融合任务状态 & \texttt{模型融合接口.py:80} \\
GET & \texttt{/fusion/task/\{task\_id\}/result} & 获取融合任务结果 & \texttt{模型融合接口.py:96} \\
POST & \texttt{/fusion/compare-strategies} & 比较不同融合策略 & \texttt{模型融合接口.py:125} \\
POST & \texttt{/fusion/optimize-weights} & 优化权重计算方法 & \texttt{模型融合接口.py:155} \\
GET & \texttt{/fusion/tasks} & 列出所有融合任务 & \texttt{模型融合接口.py:184} \\
GET & \texttt{/fusion/strategies} & 列出所有可用的融合策略 & \texttt{模型融合接口.py:196} \\
GET & \texttt{/fusion/weight-methods} & 列出所有可用的权重计算方法 & \texttt{模型融合接口.py:219} \\
GET & \texttt{/fusion/status} & 获取模型融合系统状态 & \texttt{模型融合接口.py:249} \\
\hline
\end{longtable}

\subsection{接口测试用例}

\subsubsection{创建融合任务}

\begin{lstlisting}[language=Python, caption=创建融合任务测试]
def test_create_fusion_task():
    """测试创建融合任务"""
    response = client.post("/fusion/create-task", json={
        "models": [
            {
                "model_id": "model_1",
                "model_name": "普通克里金",
                "predictions": [10.5, 11.2, 10.8],
                "variances": [0.5, 0.6, 0.4]
            },
            {
                "model_id": "model_2",
                "model_name": "泛克里金",
                "predictions": [10.8, 11.0, 10.9],
                "variances": [0.4, 0.5, 0.3]
            }
        ],
        "config": {
            "strategy": "weighted_average",
            "weight_method": "rmse_based"
        },
        "true_values": [10.6, 11.0, 10.9]
    })
    assert response.status_code == 200
    data = response.json()
    assert "task_id" in data["result"]
\end{lstlisting}

\section{路径规划接口}

\subsection{接口列表}

\begin{longtable}{|l|l|l|l|}
\hline
\textbf{方法} & \textbf{路径} & \textbf{功能} & \textbf{文件位置} \\
\hline
\endhead
POST & \texttt{/api/route-planning/plan} & 执行路径规划 & \texttt{路径规划接口.py:38} \\
POST & \texttt{/api/route-planning/optimize} & 优化现有路径 & \texttt{路径规划接口.py:52} \\
POST & \texttt{/api/route-planning/add-point} & 向路径中添加新采样点 & \texttt{路径规划接口.py:108} \\
POST & \texttt{/api/route-planning/remove-point} & 从路径中移除采样点 & \texttt{路径规划接口.py:167} \\
POST & \texttt{/api/route-planning/reorder} & 重新排序路径中的采样点 & \texttt{路径规划接口.py:225} \\
POST & \texttt{/api/route-planning/templates} & 创建路径模板 & \texttt{路径规划接口.py:283} \\
GET & \texttt{/api/route-planning/templates} & 获取所有路径模板 & \texttt{路径规划接口.py:295} \\
GET & \texttt{/api/route-planning/templates/\{template\_id\}} & 获取指定路径模板 & \texttt{路径规划接口.py:306} \\
DELETE & \texttt{/api/route-planning/templates/\{template\_id\}} & 删除路径模板 & \texttt{路径规划接口.py:321} \\
GET & \texttt{/api/route-planning/algorithms} & 获取可用的路径规划算法 & \texttt{路径规划接口.py:333} \\
GET & \texttt{/api/route-planning/vehicle-types} & 获取支持的车辆类型 & \texttt{路径规划接口.py:361} \\
GET & \texttt{/api/route-planning/optimization-goals} & 获取支持的优化目标 & \texttt{路径规划接口.py:381} \\
POST & \texttt{/api/route-planning/validate} & 验证路径是否满足约束条件 & \texttt{路径规划接口.py:401} \\
\hline
\end{longtable}

\section{API状态码定义}

\begin{table}[h]
\centering
\caption{API状态码定义}
\begin{tabular}{|l|l|l|}
\hline
\textbf{状态码} & \textbf{含义} & \textbf{使用场景} \\
\hline
200 & 成功 & 请求成功处理并返回结果 \\
201 & 创建成功 & 资源创建成功（如任务、项目） \\
400 & 请求错误 & 参数验证失败、数据格式错误 \\
401 & 未授权 & 需要认证但未提供 \\
403 & 禁止访问 & 权限不足 \\
404 & 未找到 & 资源不存在（任务、数据等） \\
409 & 冲突 & 数据冲突、重复提交 \\
422 & 数据验证失败 & 请求数据不符合验证规则 \\
429 & 请求过多 & 超出请求频率限制 \\
500 & 服务器错误 & 服务器内部错误 \\
502 & 网关错误 & 网关或代理服务器错误 \\
503 & 服务不可用 & 服务暂时不可用 \\
504 & 超时 & 请求超时 \\
\hline
\end{tabular}
\end{table}

% 第七章：测试文件与用例
\chapter{测试文件与用例}

\section{前端测试文件详细列表}

\subsection{API相关测试文件}

\begin{table}[h]
\centering
\caption{API相关测试文件}
\begin{tabular}{|l|l|l|}
\hline
\textbf{文件名} & \textbf{测试内容} & \textbf{测试用例数} \\
\hline
\texttt{tests/api.test.js} & API基础功能 & 15 \\
\texttt{tests/api/api-service.test.ts} & API服务详细测试 & 25 \\
\texttt{tests/anomaly-detect-api.test.js} & 异常检测API & 10 \\
\texttt{tests/config-api.test.js} & 配置API & 8 \\
\texttt{tests/decision-thresholds-api.test.js} & 决策阈值API & 12 \\
\texttt{tests/error-predict-api.test.js} & 误差预测API & 10 \\
\texttt{tests/model-evaluation-api.test.js} & 模型评估API & 15 \\
\texttt{tests/risk-calculate-api.test.js} & 风险计算API & 10 \\
\texttt{tests/risk-report-api.test.js} & 风险报告API & 8 \\
\texttt{tests/uncertainty-classify-api.test.js} & 不确定性分类API & 12 \\
\hline
\end{tabular}
\end{table}

\subsection{功能模块测试文件}

\begin{table}[h]
\centering
\caption{功能模块测试文件}
\begin{tabular}{|l|l|l|}
\hline
\textbf{文件名} & \textbf{测试内容} & \textbf{测试用例数} \\
\hline
\texttt{tests/advancedfilter.test.js} & 高级过滤器 & 15 \\
\texttt{tests/batchoperations.test.js} & 批量操作 & 20 \\
\texttt{tests/confirmdialog.test.js} & 确认对话框 & 10 \\
\texttt{tests/datacomparison.test.js} & 数据对比 & 12 \\
\texttt{tests/exportenhancer.test.js} & 导出增强 & 15 \\
\texttt{tests/feedbackcollector.test.js} & 反馈收集器 & 8 \\
\texttt{tests/formvalidator.test.js} & 表单验证器 & 20 \\
\texttt{tests/historymanager.test.js} & 历史管理器 & 15 \\
\texttt{tests/offlinemanager.test.js} & 离线管理器 & 18 \\
\texttt{tests/preferencespanel.test.js} & 偏好设置面板 & 10 \\
\texttt{tests/templatedownloader.test.js} & 模板下载器 & 8 \\
\texttt{tests/thememanager.test.js} & 主题管理器 & 12 \\
\hline
\end{tabular}
\end{table}

\subsection{性能和安全测试文件}

\begin{table}[h]
\centering
\caption{性能和安全测试文件}
\begin{tabular}{|l|l|l|}
\hline
\textbf{文件名} & \textbf{测试内容} & \textbf{测试用例数} \\
\hline
\texttt{tests/performance-monitor.test.ts} & 性能监控 & 20 \\
\texttt{tests/performancemonitor.test.js} & 性能监控器 & 15 \\
\texttt{tests/cache-performance.test.js} & 缓存性能 & 10 \\
\texttt{tests/security.test.ts} & 安全测试套件 & 25 \\
\texttt{tests/performance-benchmark.js} & 性能基准测试 & 8 \\
\hline
\end{tabular}
\end{table}

\section{后端测试文件详细列表}

\subsection{核心模块测试文件}

\begin{table}[h]
\centering
\caption{核心模块测试文件}
\begin{tabular}{|l|l|l|}
\hline
\textbf{文件名} & \textbf{测试内容} & \textbf{测试用例数} \\
\hline
\texttt{backend/tests/test\_model\_fusion.py} & 模型融合系统 & 30 \\
\texttt{backend/tests/test\_route\_planning.py} & 路径规划系统 & 25 \\
\texttt{backend/tests/test\_anomaly\_detector.py} & 异常检测器 & 15 \\
\texttt{backend/tests/test\_decision\_threshold\_analyzer.py} & 决策阈值分析器 & 12 \\
\texttt{backend/tests/test\_error\_predictor.py} & 误差预测器 & 15 \\
\texttt{backend/tests/test\_model\_evaluator.py} & 模型评估器 & 20 \\
\texttt{backend/tests/test\_risk\_index\_calculator.py} & 风险指数计算器 & 12 \\
\texttt{backend/tests/test\_sampling\_algorithms.py} & 采样算法 & 18 \\
\texttt{backend/tests/test\_spatial\_risk\_reporter.py} & 空间风险报告器 & 10 \\
\texttt{backend/tests/test\_uncertainty\_classifier.py} & 不确定性分类器 & 15 \\
\texttt{backend/tests/test\_websocket.py} & WebSocket服务 & 20 \\
\texttt{backend/test\_backend.py} & 后端集成测试 & 25 \\
\hline
\end{tabular}
\end{table}

\subsection{实时插值系统测试文件}

\begin{table}[h]
\centering
\caption{实时插值系统测试文件}
\begin{tabular}{|l|l|l|}
\hline
\textbf{文件名} & \textbf{测试内容} & \textbf{测试用例数} \\
\hline
\texttt{realtime\_interpolation/tests/test\_accuracy.py} & 准确性验证 & 20 \\
\texttt{realtime\_interpolation/tests/test\_performance.py} & 性能测试 & 15 \\
\texttt{realtime\_interpolation/tests/test\_integration.py} & 集成测试 & 25 \\
\texttt{realtime\_interpolation/tests/test\_cache\_system.py} & 缓存系统 & 18 \\
\texttt{realtime\_interpolation/tests/test\_incremental\_kriging.py} & 增量克里金 & 22 \\
\hline
\end{tabular}
\end{table}

\subsection{多目标优化系统测试文件}

\begin{table}[h]
\centering
\caption{多目标优化系统测试文件}
\begin{tabular}{|l|l|l|}
\hline
\textbf{文件名} & \textbf{测试内容} & \textbf{测试用例数} \\
\hline
\texttt{multi\_objective\_optimization/tests/test\_core.py} & 核心算法测试 & 35 \\
\hline
\end{tabular}
\end{table}

\section{测试用例统计}

\begin{table}[h]
\centering
\caption{测试用例统计}
\begin{tabular}{|l|l|l|l|}
\hline
\textbf{类别} & \textbf{文件数} & \textbf{测试用例数} & \textbf{占比} \\
\hline
前端单元测试 & 39 & 400 & 45\% \\
前端E2E测试 & 1 & 10 & 1\% \\
后端单元测试 & 18 & 300 & 34\% \\
后端集成测试 & 2 & 50 & 6\% \\
专业系统测试 & 6 & 120 & 14\% \\
\hline
\textbf{总计} & \textbf{66} & \textbf{880} & \textbf{100\%} \\
\hline
\end{tabular}
\end{table}

% 第八章：测试执行与报告
\chapter{测试执行与报告}

\section{测试执行流程}

\subsection{测试准备}

\begin{enumerate}
    \item 安装测试依赖
    \item 配置测试环境
    \item 准备测试数据
    \item 启动测试服务器
\end{enumerate}

\subsection{运行测试}

\subsubsection{前端测试}

\begin{lstlisting}[language=bash, caption=运行前端测试]
# 1. 安装依赖
npm install

# 2. 运行单元测试
npm test

# 3. 运行E2E测试
npm run test:e2e

# 4. 生成覆盖率报告
npm run test:coverage
\end{lstlisting}

\subsubsection{后端测试}

\begin{lstlisting}[language=bash, caption=运行后端测试]
# 1. 安装依赖
pip install -r requirements.txt

# 2. 运行单元测试
pytest

# 3. 运行特定模块测试
pytest backend/tests/test_model_fusion.py

# 4. 生成覆盖率报告
pytest --cov=app --cov-report=html
\end{lstlisting}

\subsubsection{API集成测试}

\begin{lstlisting}[language=bash, caption=运行API集成测试]
# 1. 启动后端服务器
python backend/run.py

# 2. 运行基础API测试
python api_test.py

# 3. 运行完整工作流测试
python api_test_complete.py
\end{lstlisting}

\section{测试报告}

\subsection{覆盖率报告}

\subsubsection{前端覆盖率报告}

前端覆盖率报告生成在 \texttt{coverage/} 目录下，包括：

\begin{itemize}
    \item \texttt{coverage/index.html} - HTML格式报告
    \item \texttt{coverage/coverage-final.json} - JSON格式报告
    \item \texttt{coverage/lcov.info} - LCOV格式报告
\end{itemize}

\subsubsection{后端覆盖率报告}

后端覆盖率报告生成在 \texttt{htmlcov/} 目录下，包括：

\begin{itemize}
    \item \texttt{htmlcov/index.html} - HTML格式报告
    \item 终端输出显示覆盖率摘要
\end{itemize}

\subsection{性能报告}

\subsubsection{前端性能报告}

前端性能测试结果包括：

\begin{itemize}
    \item Core Web Vitals指标（LCP、FID、CLS）
    \item 内存使用情况
    \item 资源加载时间
    \item 性能阈值检查结果
\end{itemize}

\subsubsection{后端性能报告}

后端性能测试结果包括：

\begin{itemize}
    \item API响应时间
    \item 吞吐量（QPS）
    \item 资源使用情况（CPU、内存）
    \item 算法执行时间
\end{itemize}

\subsection{测试结果分析}

\begin{enumerate}
    \item 检查测试通过率
    \item 分析失败的测试用例
    \item 检查覆盖率是否达标
    \item 识别性能瓶颈
    \item 提出改进建议
\end{enumerate}

\section{持续集成}

\subsection{CI/CD配置}

\begin{itemize}
    \item GitHub Actions自动运行测试
    \item 测试失败时阻止合并
    \item 自动生成测试报告
    \item 覆盖率阈值检查
\end{itemize}

\subsection{测试策略}

\begin{itemize}
    \item 每次提交触发单元测试
    \item 每次PR触发完整测试套件
    \item 每日运行性能测试
    \item 每周运行E2E测试
\end{itemize}

\section{测试最佳实践}

\subsection{编写测试用例}

\begin{itemize}
    \item 遵循AAA模式（Arrange-Act-Assert）
    \item 测试应该独立且可重复
    \item 使用有意义的测试名称
    \item 测试正常情况和边界情况
    \item 避免测试实现细节
\end{itemize}

\subsection{Mock和Stub}

\begin{itemize}
    \item Mock外部依赖
    \item Stub函数返回值
    \item 隔离被测单元
    \item 避免过度Mock
\end{itemize}

\subsection{测试数据管理}

\begin{itemize}
    \item 使用测试数据工厂
    \item 避免硬编码测试数据
    \item 使用测试数据库
    \item 清理测试数据
\end{itemize}

\section{测试维护}

\subsection{定期审查}

\begin{itemize}
    \item 每月审查测试覆盖率
    \item 每季度审查测试策略
    \item 定期更新测试依赖
    \item 优化慢速测试
\end{itemize}

\subsection{问题追踪}

\begin{itemize}
    \item 记录测试失败原因
    \item 追踪测试问题修复
    \item 定期review测试代码
    \item 重构重复测试代码
\end{itemize}

% 附录：数据格式与参考文档
\appendix

\chapter{数据格式参考}

\section{实时插值数据格式}

\subsection{DataPoint数据结构}

\begin{lstlisting}[language=Python, caption=DataPoint数据结构]
@dataclass
class DataPoint:
    """数据点"""
    x: float              # X坐标
    y: float              # Y坐标
    value: float          # 观测值
    id: str               # 唯一标识符
    timestamp: Optional[datetime] = None   # 时间戳
    metadata: Optional[Dict[str, Any]] = None  # 元数据
\end{lstlisting}

\subsection{Subscription数据结构}

\begin{lstlisting}[language=Python, caption=Subscription数据结构]
@dataclass
class Subscription:
    """订阅"""
    subscription_id: str
    data_type: str
    spatial_extent: BoundingBox
    update_frequency: int
    interpolation_params: VariogramModel
    notification_config: Dict[str, Any]
    created_at: datetime
    updated_at: datetime
\end{lstlisting}

\section{多目标优化数据格式}

\subsection{OptimizationRequest数据结构}

\begin{lstlisting}[language=Python, caption=OptimizationRequest数据结构]
@dataclass
class OptimizationRequest:
    """优化请求"""
    variance_grid: Dict[str, Any]
    existing_points: List[Dict[str, float]]
    n_samples: int
    weights: Dict[str, float]
    constraints: Dict[str, Any]
    algorithm: str
    algorithm_params: Dict[str, Any]
    is_async: bool
\end{lstlisting}

\subsection{ParetoSolution数据结构}

\begin{lstlisting}[language=Python, caption=ParetoSolution数据结构]
@dataclass
class ParetoSolution:
    """帕累托解"""
    id: int
    objectives: Dict[str, float]
    sampling_points: List[Dict[str, Any]]
    metrics: Dict[str, float]
    rank: int
    crowding_distance: float
\end{lstlisting}

\section{模型融合数据格式}

\subsection{ModelInput数据结构}

\begin{lstlisting}[language=Python, caption=ModelInput数据结构]
@dataclass
class ModelInput:
    """模型输入"""
    model_id: str
    model_name: str
    predictions: List[float]
    variances: List[float]
\end{lstlisting}

\subsection{FusionResult数据结构}

\begin{lstlisting}[language=Python, caption=FusionResult数据结构]
@dataclass
class FusionResult:
    """融合结果"""
    task_id: str
    fused_predictions: List[float]
    fused_variances: List[float]
    weights: Dict[str, float]
    metrics: Dict[str, float]
    strategy: str
    weight_method: str
\end{lstlisting}

\chapter{参考文档}

\section{项目文档}

\begin{itemize}
    \item \texttt{docs/系统架构.md} - 系统架构文档
    \item \texttt{docs/开发指南.md} - 开发指南
    \item \texttt{docs/使用手册.md} - 用户手册
    \item \texttt{docs/插件开发.md} - 插件开发文档
\end{itemize}

\section{API文档}

\begin{itemize}
    \item \texttt{realtime\_interpolation/接口设计文档.md} - 实时插值API文档
    \item \texttt{multi\_objective\_optimization/API接口设计文档.md} - 多目标优化API文档
    \item \texttt{docs/插件架构.md} - 插件API文档
\end{itemize}

\section{算法文档}

\begin{itemize}
    \item \texttt{realtime\_interpolation/算法设计文档.md} - 实时插值算法文档
    \item \texttt{multi\_objective\_optimization/算法设计文档.md} - 多目标优化算法文档
    \item \texttt{docs/模型融合系统文档.md} - 模型融合算法文档
\end{itemize}

\section{外部资源}

\begin{itemize}
    \item Vitest官方文档：\url{https://vitest.dev/}
    \item Playwright官方文档：\url{https://playwright.dev/}
    \item pytest官方文档：\url{https://docs.pytest.org/}
    \item FastAPI官方文档：\url{https://fastapi.tiangolo.com/}
\end{itemize}

\chapter{术语表}

\begin{description}
    \item[克里金插值] 一种基于统计学的空间插值方法
    \item[NSGA-II] 非支配排序遗传算法II，用于多目标优化
    \item[帕累托前沿] 多目标优化中的最优解集合
    \item[WebSocket] 一种在单个TCP连接上进行全双工通讯的协议
    \item[E2E测试] 端到端测试，测试完整的用户工作流
    \item[覆盖率] 测试覆盖的代码比例
    \item[Core Web Vitals] Google提出的Web性能指标
    \item[LCP] Largest Contentful Paint，最大内容绘制时间
    \item[FID] First Input Delay，首次输入延迟
    \item[CLS] Cumulative Layout Shift，累积布局偏移
    \item[API] Application Programming Interface，应用程序接口
    \item[REST] Representational State Transfer，表述性状态转移
    \item[JSON] JavaScript Object Notation，一种轻量级的数据交换格式
    \item[GeoJSON] 一种基于JSON的地理空间数据格式
    \item[WebSocket] 一种全双工通信协议
    \item[SSE] Server-Sent Events，服务器推送事件
    \item[Mock] 模拟对象，用于隔离测试
    \item[Stub] 桩函数，用于提供预设的返回值
\end{description}

\end{document}